% =====================================================================================================
% Chương 2 — bản viết lại, phiên bản 2 (2026-10-12). Dùng như \input{noi-dung/06-chuong2} của main.tex.
% Bản 1 (06-chuong2.tex) giữ nguyên; bản 2 rút gọn khoảng 15% văn xuôi cho vừa 6,5 trang và sửa sáu chỗ diễn đạt
% so với bài gốc (06-chuong2-NOTES.md mục 3). Tài liệu tham khảo: 06-chuong2-BIB.md.
% Mỗi dòng "% EV:" trỏ tới EVIDENCE-REGISTER.md (A01…K52), các phụ lục 1–2 (N-xx, X3-xx) hoặc
% EVIDENCE-REGISTER-ADDENDUM-5-CH2.md (C2-01…C2-24, các dòng kiểm nguồn mới của chương này).
% Số trang trong EV là trang của bài; trang PDF ghi riêng khi khác.
% =====================================================================================================

\chapter{Cơ sở lý thuyết và các công trình liên quan}
\label{chuong:coso}

Chương này trình bày bài toán ghép chuyến động mà khóa luận dựa vào, các cách nghiên cứu trước đã đo và ràng buộc
lời hứa với khách, và những gì đã biết về dự báo thời gian đi và giá trị của độ tin cậy. Mỗi mục tách phần là lý
thuyết hay kết quả của người khác khỏi phần RideBound dùng lại; định nghĩa hình thức của luật C và các luật đối chứng
nằm ở Chương~\ref{chuong:phuongphap}.

% -----------------------------------------------------------------------------------------------------
\section{Bài toán ghép chuyến động và kiến trúc yêu cầu--chuyến--xe}
\label{muc:c2-baitoan}

Trong dịch vụ đi chung xe, mỗi yêu cầu gồm điểm đón, điểm trả và thời điểm đặt; một đội xe có sức chứa giới hạn chở
nhiều khách cùng lúc. Bộ điều phối (dispatcher) quyết định nhận hay từ chối từng yêu cầu, giao nó cho xe nào và xếp
các điểm đón, điểm trả theo thứ tự nào. Bài toán này thuộc họ bài toán đưa đón theo yêu cầu (Dial-a-Ride Problem,
DARP): thiết kế các tuyến xe có chi phí nhỏ nhất để chở khách chung theo nhu cầu, trong khi vẫn giữ các yêu cầu về mức
dịch vụ~\cite{molenbruch2017}. Ở dạng động, chi tiết của yêu cầu chỉ được biết khi khách gửi~\cite{schulz2026}, nên
bộ điều phối giải lại bài toán nhiều lần trong ngày với trạng thái lúc đó; mỗi lần giải là một lượt quyết định
(decision epoch).
% EV: J25; C2-12 (Molenbruch, PDF tr. 3: định nghĩa DARP); C2-07 (Schulz–Pfeiffer tr. 549: yêu cầu không biết trước); X3-01

Alonso-Mora và cộng sự~\cite{alonsomora2017} đề xuất một cách giải theo lô cho xe đi chung có sức chứa lớn. Yêu cầu
được gom trong một khoảng ngắn, 30 giây trong thí nghiệm của họ, rồi được gán cùng một lần. Bước đầu dựng một đồ thị
nối hai yêu cầu, hay một yêu cầu với một xe, khi chúng chở chung được mà vẫn thỏa các ràng buộc dịch vụ. Từ đó, bước
sau dựng đồ thị yêu cầu--chuyến--xe (request--trip--vehicle graph), trong đó một chuyến là một nhóm yêu cầu mà một xe
chở được. Một nhóm chỉ cần kiểm khi mọi nhóm con bớt đi một yêu cầu đều khả thi với cùng xe, nên các nhóm được dựng
dần theo kích thước. Hörl và Zwick~\cite{horl2022} cài thuật toán này trong bộ mô phỏng MATSim, với các lượt quyết
định thường cách nhau 30 giây; FleetPy cũng có thuật toán này~\cite{engelhardt2026}.
% EV: J17; C2-10 (Alonso-Mora, bản PNAS sớm tr. 3: đồ thị yêu cầu–xe, yêu cầu–chuyến–xe, kiểm nhóm con; tr. 4: lô 30 s); C2-02 (Hörl–Zwick tr. 4, PDF tr. 5); J18; C2-11 (FleetPy tr. 11: thuật toán Alonso-Mora)

Bước chọn là một quy hoạch nguyên (integer linear program). Gọi $\epsilon_{ij} \in \{0, 1\}$ bằng 1 khi xe $j$ được
giao chuyến $T_i$, và $\chi_k \in \{0, 1\}$ bằng 1 khi yêu cầu $k$ không được xe nào phục vụ. Bài toán là
\begin{equation}
\min \sum_{(i,j)} c_{ij}\,\epsilon_{ij} + \sum_{k} c_{\mathrm{ko}}\,\chi_k
\quad\text{với}\quad
\sum_{i} \epsilon_{ij} \le 1 \;\;\text{với mọi xe } j, \qquad
\sum_{(i,j):\, k \in T_i} \epsilon_{ij} + \chi_k = 1 \;\;\text{với mọi yêu cầu } k .
\label{eq:c2-chon}
\end{equation}
Trong đó, tổng thứ nhất lấy trên các cặp chuyến--xe có trong đồ thị; $c_{ij}$ là tổng độ trễ của mọi khách khi xe $j$
chạy chuyến $T_i$, kể cả khách đang ngồi trên xe; $c_{\mathrm{ko}}$ là một hằng số lớn phạt mỗi yêu cầu không được
phục vụ. Độ trễ của một khách là hiệu giữa giờ trả và giờ sớm nhất khách có thể tới nếu đi thẳng ngay lúc gọi, nên
gồm cả thời gian chờ. Ràng buộc thứ nhất cho mỗi xe nhận nhiều nhất một chuyến; ràng buộc thứ hai buộc mỗi yêu cầu
hoặc thuộc đúng một chuyến được chọn, hoặc không được phục vụ.
% EV: C2-10 (Alonso-Mora tr. 3: hàm chi phí (1), định nghĩa độ trễ; tr. 4: giải thuật 1, dòng 2–4)

Hàm mục tiêu~\eqref{eq:c2-chon} đã tính thiệt hại của khách cũ, vì ghép thêm khách làm độ trễ của khách đang trên xe
tăng. Tuy vậy, đó là một số hạng mềm, cộng chung trên mọi khách và đo so với chuyến đi thẳng, không so với giờ đã báo
cho khách. Một khách vì thế có thể bị dời giờ trả nhiều lần, miễn là mỗi lần tổng chi phí nhỏ hơn. Các hàm mục tiêu
kiểu này không định giá việc sửa lời hứa; Fielbaum và Alonso-Mora đã đo và kiểm soát hiện tượng đó
(mục~\ref{muc:c2-dolech}).
% EV: C2-10 (tr. 3: chi phí gồm độ trễ của khách trên xe; độ trễ so với đi thẳng); K30; câu thứ ba là suy luận từ định nghĩa, không phải kết quả đo

RideBound giữ bước chọn theo lối này nhưng không dựng đồ thị yêu cầu--chuyến--xe. Với mỗi xe, bộ sinh chỉ chèn thêm
(insertion-only generator) tạo phương án giữ nguyên kế hoạch và các phương án chèn điểm đón, điểm trả của yêu cầu mới
vào kế hoạch cũ, không bỏ và không đảo điểm cũ. Bộ chọn giao mỗi xe đúng một phương án, mỗi yêu cầu mới thuộc nhiều
nhất một xe, và so các phương án theo một thứ tự tiêu chí thay cho một tổng chi phí; khách đã nhận không được chuyển
sang xe khác (mục~\ref{muc:baitoan-phatbieu}). Các ràng buộc dịch vụ có cùng dạng với các bài trên; giá trị dùng
trong khóa luận ở Bảng~\ref{bang:c2-dichvu}.
% EV: X3-03 (GEN:196-212, 231-246, 548-575); X3-22 (ORT:96-125); X3-07; X3-39; A14; A24

\begin{table}[htbp]
\centering
\caption{Ràng buộc dịch vụ của một phương án: cách đặt ở một số nghiên cứu trước và giá trị trong các tập thử của
khóa luận. Giá trị của khóa luận là cấu hình của tập thử, không phải chuẩn ngành; mỗi nghiên cứu tính từ một mốc
khác nên các con số không so trực tiếp được.}
\label{bang:c2-dichvu}
\small
\begin{tabular}{@{}p{2.4cm}p{3.4cm}p{5.9cm}p{3.1cm}@{}}
\toprule
\textbf{Ràng buộc} & \textbf{Ý nghĩa} & \textbf{Ở một số nghiên cứu trước} & \textbf{Trong các tập thử}\\
\midrule
Cửa sổ đón & Khách phải được đón trong một khoảng tối đa kể từ lúc gọi xe & 2, 5 hoặc 7
phút~\cite{alonsomora2017}; 5 phút~\cite{fielbaum2020,elmi2026}; 480 giây ở nghiên cứu điển hình của
FleetPy~\cite{engelhardt2026} & 600 giây kể từ lúc gọi\\
% EV: C2-10 (tr. 5: 120/300/420 s); C2-06 (FAM tr. 18, bảng 6); C2-08 (Elmi tr. 15); C2-11 (FleetPy tr. 11); B03; C2-22 (NORM:444)
Thời gian ngồi xe hay độ trễ & Giới hạn phần thời gian khách mất thêm so với đi thẳng & Độ trễ tổng, gồm cả thời gian
chờ, bằng hai lần cửa sổ đón~\cite{alonsomora2017}; độ trễ tổng 10 phút~\cite{fielbaum2020,elmi2026}; phần vượt thời
gian đi thẳng tính theo tỉ lệ của nó~\cite{molenbruch2017}; độ vòng tương đối 40\%~\cite{engelhardt2026} & Thời gian
ngồi trên xe không quá 1,5 lần thời gian đi thẳng\\
% EV: C2-10 (tr. 5: ∆ bằng hai lần Ω, gồm thời gian chờ); C2-06 (bảng 6); C2-08 (tr. 15); C2-12 (PDF tr. 5); C2-11 (tr. 11); B03; C2-22 (PhysicalPlanValidator.cs:594)
Sức chứa & Số khách tối đa cùng lúc trên một xe & 1, 2, 4 hoặc 10 chỗ~\cite{alonsomora2017}; 4
chỗ~\cite{engelhardt2026}; 6 chỗ~\cite{schulz2026} & 4 chỗ; 8 xe mỗi kịch bản\\
% EV: C2-10 (tr. 5); C2-11 (tr. 11: 180 xe 4 chỗ); C2-07 (tr. 562); B03
\bottomrule
\end{tabular}
\end{table}

Các con số trong Bảng~\ref{bang:c2-dichvu} không so trực tiếp được vì mỗi bài tính từ một mốc khác: độ trễ tổng của
Alonso-Mora và cộng sự gồm cả thời gian chờ, còn giới hạn của khóa luận chỉ tính thời gian ngồi trên xe. Các giá trị
của khóa luận chưa được hiệu chỉnh theo sở thích của khách.
% EV: C2-10 (tr. 5); C2-22 (PhysicalPlanValidator.cs:594: thời gian ngồi xe = giờ trả − giờ đón); C2-23 (D03:252, 259-260)

% -----------------------------------------------------------------------------------------------------
\section{Đo lường và kiểm soát độ bất ổn của lời hứa}
\label{muc:c2-dolech}

Fielbaum và Alonso-Mora~\cite{fielbaum2020} phân loại độ bất ổn (unreliability) của dịch vụ đi chung xe theo nguồn:
từ khách, như khách vắng mặt hay hủy chuyến, và từ hãng xe, khi cách gom và gán yêu cầu thay đổi mỗi lần có yêu cầu
mới; họ tập trung vào nguồn thứ hai. Giá trị báo ban đầu của một khách là giá trị sẽ đạt nếu lộ trình của xe không đổi
kể từ lần gán thứ nhất. Chỉ số của họ là hiệu giữa giá trị thực tế và giá trị báo ban đầu của thời gian chờ, độ vòng
và độ trễ tổng, với giá trị âm đặt bằng 0 để chỉ giữ phần bất lợi; họ còn đếm số lần giá trị dự kiến tăng lên.
% EV: J06; X3-35; C2-06 (FAM tr. 4: định nghĩa 1–2, tập trung vào hãng xe, giá trị báo ban đầu; tr. 5: các hiệu, cắt giá trị âm; tr. 7: định nghĩa 9)

Thí nghiệm của họ dùng mạng Manhattan 4.091 nút, 9.453 cạnh, 2.000 xe 4 chỗ và 19.610 yêu cầu trong khung 13--15 giờ
ngày 15/01/2013. Khung thấp điểm được chọn vì khi đó thời gian đi khá ổn định, cho phép bỏ qua độ bất ổn do tắc
đường. Khi không hạn chế việc chuyển khách sang xe khác, hơn một phần ba số yêu cầu không bị từ chối ngay gặp ít nhất
một thay đổi; khi thời gian chờ của một khách tăng, phần tăng tương đương thời gian chờ trung bình của cả hệ thống.
% EV: J07; C2-06 (tr. 9: dữ liệu và lý do chọn giờ thấp điểm; tr. 1: hơn một phần ba, khi không hạn chế chuyển xe)

Ở mục 6, họ kiểm soát độ bất ổn ngay trong quyết định. Trong kịch bản kiểm soát hoàn toàn, khách đã được gán không bị
từ chối về sau và thời gian chờ tối đa của khách được hạ về giá trị dự kiến; tỉ lệ yêu cầu bị từ chối tăng từ 34,5\%
lên 45,6\%, và chi phí trung bình của khách theo mô hình tăng từ 9,56 lên 11,6. Kịch bản trung gian chỉ cấm từ chối
khách đã nhận: tỉ lệ từ chối còn 31,2\%, nhưng thời gian chờ và độ vòng của khách được nhận tăng khoảng 50\%. Họ kết
luận có đánh đổi giữa độ tin cậy và các chỉ số chất lượng khác, và nêu câu hỏi mở là kiểm soát độ bất ổn mà không làm
tăng tỉ lệ từ chối.
% EV: J08 (số trong bảng 5); C2-06 (tr. 17: hai kịch bản, bảng 5 đọc bằng pdftotext -table; tr. 18: câu hỏi mở); J09; K32

Elmi và cộng sự~\cite{elmi2026} tiếp tục hướng này theo từng khách. Độ chùng (slack) của một khách là hiệu giữa thời
gian chờ tối đa, hay độ trễ tối đa, và giá trị dự kiến lúc gán; đó là phần thời gian còn cho phép thay đổi về sau.
Khách có điểm khả năng đi chung thấp nhận độ chùng 0, tức được bảo đảm không bị đổi giờ, còn khách khác giữ toàn bộ độ
chùng. Ở cùng tỉ lệ từ chối 29\% với cách siết đều cho mọi khách, cách này bảo đảm chuyến không đổi giờ cho 83,9\% số
khách ngay lúc gán. Bài giả định thời gian đi trên mỗi cạnh không đổi và cố ý không xét độ bất ổn do khách hay do giao
thông.
% EV: C2-08 (Elmi tr. 1: 29%, 83,9%; tr. 2: không xét nguồn từ khách và từ mạng giao thông; tr. 5: thời gian đi không đổi, công thức 2; tr. 12–13: hai cách điều khiển, độ chùng 0); J12 (nay đã mở bài gốc)

Thước đo của Fielbaum và Alonso-Mora so giá trị cuối với giá trị báo ban đầu. Vì vậy một khách bị dời muộn 5 phút rồi
được kéo về đúng giờ cũ có hiệu bằng 0, dù đã nhận hai lần thay đổi; số lần tăng thì được đếm nhưng không mang độ
lớn. Khóa luận kế thừa ý tách phần thay đổi do quyết định của hãng xe và kiểm soát nó trước khi quyết định, nhưng đo
từng lần sửa so với giờ trả khi giữ nguyên kế hoạch và cộng độ lớn các lần đó qua cả chuyến (mục~\ref{muc:bamoc}).
Trên tập kiểm tra chính, mỗi lần điều phối kéo giờ trả sớm hơn chỉ sớm 1 mili giây, nên khác biệt này nằm ở định
nghĩa, chưa thấy ở kết quả.
% EV: C2-06 (tr. 5: hiệu giữa giá trị thực tế và báo ban đầu; tr. 7: chỉ đếm số lần tăng); ví dụ 5 phút dựng tay từ định nghĩa; A01; A04; K30; K32; ADDENDUM-1 mục 0.2 (ADD1:104-111); 07-chuong3-v3.tex:287-291

Đánh đổi giữa ổn định và hiệu quả cũng đã được nghiên cứu ở các lĩnh vực khác: định tuyến nhất quán, nơi khách được
phục vụ vào giờ quen thuộc bởi cùng một tài xế~\cite{kovacs2014}, và đóng băng một phần lịch sản xuất để giảm thay
đổi lịch, đổi lại chi phí sản xuất và tồn kho có thể tăng~\cite{sridharan1987}.
% EV: J26; C2-20 (Kovacs tr. 192, tóm tắt); C2-21 (Sridharan tr. 1137, tóm tắt)

% -----------------------------------------------------------------------------------------------------
\section{Các dạng ràng buộc lời hứa}
\label{muc:c2-rangbuoc}

Ngoài cửa sổ đón và giới hạn thời gian ngồi xe, các nghiên cứu trước còn ràng buộc chính lời hứa với khách.
Bảng~\ref{bang:c2-dangrangbuoc} so các cách đó theo năm câu hỏi: so với mốc nào, phần do tắc đường có bị tính không,
áp cho khách nào, có cộng qua cả chuyến không, và có ghi riêng cho từng khách không.

Nhiều cách giữ một giờ đã báo. Alonso-Mora và cộng sự hạ giờ đón muộn nhất về giờ đón dự kiến khi yêu cầu được ghép
xe, và chỉ chuyển yêu cầu sang xe khác nếu thời gian chờ không tăng~\cite{alonsomora2017}. Lu và
Kühnel~\cite{lu2024} cho giờ đón muộn hơn giờ hẹn ban đầu nhiều nhất một khoảng đệm, gợi ý 2--3 phút để chừa cho tắc
đường và các yếu tố vận hành khác. Schulz và Pfeiffer~\cite{schulz2026} đặt hạn chót trả bằng giờ trả trong lời mời
(offer) cộng một độ chùng, mặc định 900 giây; khi quét độ chùng tới 30 phút, họ thấy 10--15 phút là đánh đổi hợp lý.
Mt-KaRRi~\cite{laupichler2026} giới hạn thời gian chuyến của khách ở 1,4 lần thời gian chuyến đã báo lúc khách chấp
nhận, cộng 10 phút. Hai bài sau dùng thời gian đi cố định.
% EV: C2-10 (Alonso-Mora tr. 4); C2-14 (Lu–Kühnel tr. 4: công thức 3, đệm 2–3 phút); J10; C2-07 (Schulz–Pfeiffer tr. 549–551, 562, 569); J32; C2-09 (Mt-KaRRi tr. 4, 5, 27)

Gần luật C nhất là chiến lược B của Hörl và Zwick~\cite{horl2022}, khi họ xét thời gian đi thay đổi trong mô phỏng
MATSim với nhu cầu thật của MOIA ở Hamburg. Ở đầu mỗi lượt, họ duyệt lại tuyến đang gán theo thời gian đi hiện tại để
có giờ trả cập nhật. Theo công thức 7 của bài, giờ trả bắt buộc của các yêu cầu trên xe đang xét bằng giá trị lớn hơn
giữa giờ trả cập nhật và giờ trả muộn nhất; tuyến phải thỏa giờ bắt buộc như ràng buộc cứng (ràng buộc 9--10), còn
mỗi giây vượt giờ trả muộn nhất gốc bị phạt trong hàm mục tiêu (công thức 11). Nhờ vậy xe vẫn được đón thêm khách, miễn
là khách trên xe không bị trả muộn hơn mức sớm nhất có thể với tắc đường hiện tại, và thuật toán không rơi vào tình
huống không còn thứ tự trả khách khả thi. Giờ đón được xử lý tương tự (chiến lược C). Họ chỉ ra rằng khi thời gian đi
bị ước lượng thấp, hãng xe phải chọn giữa từ chối khách đã nhận và trả khách muộn hơn giờ đã báo; trong kịch bản có
thời gian đi thực tế, tác động này nhỏ.
% EV: J01; J02; J03; J04; X3-36; C2-04 (Hörl–Zwick tr. 7–8, PDF tr. 8–9: chiến lược B, C, công thức 5–11); C2-05 (tr. 14, PDF tr. 15); K26 (không lấy cứng/mềm làm khác biệt)

\begin{table}[htbp]
\centering
\caption{Các dạng ràng buộc lời hứa ở nghiên cứu trước và luật C, mô tả bằng lời. "Không" ở cột "Cộng qua cả
chuyến?" nghĩa là ràng buộc so một giờ với một mốc, không cộng độ lớn các lần sửa. Định nghĩa hình thức của luật C và
các luật đối chứng ở Chương~\ref{chuong:phuongphap}.}
\label{bang:c2-dangrangbuoc}
\footnotesize
\setlength{\tabcolsep}{3pt}
\begin{tabular}{@{}p{2.7cm}p{2.6cm}p{2.5cm}p{2.5cm}p{2.3cm}p{2.3cm}@{}}
\toprule
\textbf{Dạng ràng buộc} & \textbf{So với mốc nào} & \textbf{Phần do tắc đường có bị tính?} & \textbf{Áp cho ai} &
\textbf{Cộng qua cả chuyến?} & \textbf{Ghi riêng cho từng khách?}\\
\midrule
Cửa sổ đón~\cite{alonsomora2017} & Giờ gọi xe & Có & Mọi yêu cầu, tới lúc lên xe & Không & Không\\
% EV: C2-10 (tr. 3); B03
Thời gian ngồi xe hay độ trễ tổng~\cite{alonsomora2017,molenbruch2017} & Thời gian đi thẳng & Có & Mọi khách & Không
& Không\\
% EV: C2-10 (tr. 3); C2-12 (PDF tr. 5)
Giữ giờ đón đã báo~\cite{alonsomora2017,lu2024} & Giờ đón dự kiến lúc ghép xe, hay giờ đón hẹn ban đầu cộng 2--3
phút & Có; khoảng đệm chừa cho tắc đường & Khách chưa lên xe & Không & Không\\
% EV: C2-10 (tr. 4); C2-14 (tr. 4)
Hạn chót từ lời mời đầu~\cite{schulz2026} & Giờ trả trong lời mời, cộng độ chùng 900 giây & Bài dùng thời gian đi cho
trước; khi thời gian đi đổi thì có & Mọi khách đã nhận lời mời & Không: so giờ trả mới với lời mời & Không; mỗi khách
một hạn chót\\
% EV: J10; C2-07 (tr. 549–551, 562); A10 (riêng tắc đường có thể làm hạn chót bị vượt)
Giờ trả cập nhật và hạn trả được nới~\cite{horl2022} (công thức 7, ràng buộc 9--10, phạt 11) & Giờ trả muộn nhất;
giờ cập nhật tính lại ở mỗi lượt & Hạn được nới tới giờ cập nhật, nên riêng tắc đường không làm tuyến bị loại & Khách
của xe đang xét; có tác dụng khi giờ cập nhật đã vượt hạn & Không & Không; phần vượt hạn gốc bị phạt\\
% EV: J01; J03; J05; C2-04 (tr. 7–8)
Bảo đảm theo từng khách~\cite{elmi2026} & Thời gian chờ và độ trễ dự kiến lúc gán & Ngoài phạm vi: thời gian đi trên
mỗi cạnh không đổi & Khách có điểm đi chung thấp không bị đổi giờ; khách khác trong độ chùng & Không: so với giá trị
lúc gán & Không; chỉ chọn khách theo điểm\\
% EV: C2-08 (tr. 2, 5, 12–13)
Giới hạn thời gian chuyến của Mt-KaRRi~\cite{laupichler2026} & 1,4 lần thời gian chuyến đã báo lúc chấp nhận, cộng
10 phút & Thời gian đi trên mỗi cạnh cố định & Mọi khách đã chấp nhận & Không: mốc không đổi sau mỗi lần sửa & Không\\
% EV: J32; C2-09 (tr. 4, 5, 27)
\textbf{Luật C (khóa luận)} & Giờ trả nếu xe giữ nguyên kế hoạch, tính lại ở mỗi lượt & Không: chỉ tính phần dời do
điều phối & Mọi khách đã xác nhận đặt xe & Có: cộng độ lớn từng lần dời & Có: sổ theo từng khách\\
% EV: A01; A02; A06; A13; K43; X3-25
\bottomrule
\end{tabular}
\end{table}

Luật C dùng lại mốc giờ trả cập nhật của Hörl và Zwick nhưng giới hạn một đại lượng khác: phần dời do điều phối so với
mốc đó, cộng độ lớn qua cả chuyến, với cùng một trần cho mọi khách đã xác nhận, không nới theo tắc đường, và ghi vào sổ
theo từng khách (mục~\ref{muc:luatc}). Nếu cộng có dấu cả phần dời do điều phối lẫn phần thay đổi khi giữ nguyên kế
hoạch thì được lại hạn chót tính từ lời hứa đầu; nếu cộng độ lớn phần khách thấy thì được luật trần độ dời khách thấy.
Hai luật này làm luật đối chứng (mục~\ref{muc:doichung}). Bảo vệ từng khách, giới hạn giờ đến, chèn khách động và
chuyển khách sang xe khác đều đã có trong các nghiên cứu trên, nên khóa luận không nhận các ý đó là mới.
% EV: K01; J05; A12; X3-36; J38; J39; K47; K46 (không dùng "duy nhất", "đầu tiên"); C2-07, C2-08, C2-09, C2-10 (chèn động, bảo vệ từng khách, giới hạn giờ đến, chuyển xe)

% -----------------------------------------------------------------------------------------------------
\section{Dự báo thời gian đi và lời hứa theo dự báo}
\label{muc:c2-dubao}

Bộ điều phối thường tính giờ đón, giờ trả bằng thời gian đi lúc ra quyết định. Khi đường đông dần, giờ tính theo cách
này sớm hơn giờ xe thực sự tới, dù xe giữ nguyên kế hoạch; đó là tình huống thời gian đi bị ước lượng thấp mà Hörl và
Zwick đã xét, trong mô phỏng có tốc độ cập nhật mỗi 60 phút từ dữ liệu định vị của TomTom~\cite{horl2022}.
% EV: A05; F10; C2-05 (Hörl–Zwick tr. 10, PDF tr. 11: tốc độ cập nhật mỗi 60 phút; tr. 14: ước lượng thấp)

Dự báo giờ đến từ dữ liệu là thực hành của các nền tảng gọi xe. Uber dùng giờ đến dự kiến để tính giá cước, ước lượng
giờ đón và giờ trả, và ghép khách với tài xế; mô hình học sâu DeeprETA hiệu chỉnh giờ đến do bộ tìm đường tính, với
các đặc trưng gồm phút trong ngày, ngày trong tuần, tốc độ thời gian thực và tốc độ lịch sử. Nhóm tác giả cũng lưu ý
giờ đến hiển thị cho khách có thể cần một phân vị thay cho giá trị trung bình~\cite{hu2022}. Với dịch vụ đáp ứng theo
nhu cầu, Leonardi và cộng sự~\cite{leonardi2024} học từ một số chuyến đã quan sát trong mô phỏng một khoảng thời gian
chuyến đúng trong hơn 90\% số lần, và đề xuất lấy mức dự báo được làm một tiêu chí khi định cỡ dịch vụ. Vì vậy lời hứa
theo dự báo (forecast-aware promise), tức lời hứa đã cộng sẵn mức tăng tắc đường dự báo, không phải ý mới của khóa
luận.
% EV: J35 (nay có nguồn sơ cấp C2-15); C2-15 (DeeprETA tr. 1: công dụng, phân vị; tr. 3: bảng 1); J21; C2-13 (Leonardi PDF tr. 2: tóm tắt, mô phỏng VISUM, hơn 90%; tr. 4: định cỡ theo mức dự báo được); K37

Một dự báo dùng cho lời hứa chỉ được dùng thông tin có trước lúc hứa: hồ sơ tắc đường của các ngày đã qua và mức tắc
đường đo được lúc đó. Dữ liệu chuyến công khai của Ủy ban Taxi và Xe limousine New York (Taxi \& Limousine
Commission, TLC) được công bố theo tháng, thường trễ khoảng hai tháng~\cite{tlc}. Hồ sơ lập từ các ngày ngay trước, như
ở mục~\ref{muc:dubao}, vì vậy giả định hãng xe có sẵn dữ liệu chuyến của chính mình.
% EV: B27 (nay VERIFIED-by-me, C2-16: trang TLC đọc ngày 2026-10-12); F02; câu cuối là suy luận của người viết, ghi trong NOTES

Bộ dữ liệu chuẩn Manhattan của FleetPy~\cite{engelhardt2026} lấy nhu cầu từ chuyến taxi tuần 11--18/11/2018. Để hiệu
chỉnh thời gian đi, bài so thời lượng chuyến được ghi với thời gian đi đường nhanh nhất theo tốc độ tối đa cho phép,
và tính cho mỗi giờ trong ngày một hệ số nhân; gói dữ liệu còn có tệp hệ số theo khung 15 phút cho từng ngày. Khóa
luận dùng tệp 15 phút này làm hệ số tắc đường (travel-time factor) chung cho cả thành phố, không phải tốc độ của từng
đoạn đường (mục~\ref{muc:c4-kichban}). Vì cả thành phố dùng một hệ số và mô phỏng không có nhiễu ngẫu nhiên, thế giới
mô phỏng gần như tất định; kết quả của lời hứa theo dự báo ở Chương~\ref{chuong:thucnghiem} nên được đọc như mức trần
cho loại dự báo theo giờ trong ngày, không phải hiệu quả ngoài đời.
% EV: J18; C2-11 (FleetPy tr. 10: tuần dữ liệu, hệ số theo giờ; tệp *_tt_factors.csv 96 dòng cách 900 s; PREF:51-55, 67-93); B06; I02; F10

% -----------------------------------------------------------------------------------------------------
\section{Độ tin cậy có giá trị gì với khách}
\label{muc:c2-giatri}

Các nghiên cứu về giá trị thời gian cho thấy khách coi việc tới muộn nặng hơn thời gian ngồi xe. Phân tích tổng hợp
của Wardman và cộng sự~\cite{wardman2026} cho mỗi phút tới đích muộn so với giờ biểu hay giờ dự kiến một hệ số trung
bình 2,53 lần phút ngồi xe (15 quan sát, 2 nghiên cứu) hoặc 3,12 lần (66 quan sát, 27 nghiên cứu), tùy cách hỏi. Với
dịch vụ đi chung xe, Alonso-González và cộng sự~\cite{alonsogonzalez2020} khảo sát 1.006 người ở Hà Lan và thấy tỉ số
giữa giá trị độ tin cậy và giá trị thời gian (reliability ratio) khoảng 0,5; một khoản trễ nhỏ không báo trước được
định giá thấp hơn thời gian đã báo, và người trả lời không được đặt vào tình huống tới muộn. Trong một thí nghiệm thực
địa trên 1.425.745 chuyến, Yu và cộng sự~\cite{yu2022} thấy cả độ lớn của thời gian chờ báo ban đầu lẫn số lần cập
nhật đều ảnh hưởng tới việc khách bỏ yêu cầu.\footnote{Với bài này, khóa luận chỉ đọc phần tóm tắt.}
% EV: J23; C2-17 (Wardman tr. 2: hệ số so với thời gian ngồi xe; tr. 3: định nghĩa; tr. 10: bảng 5); J24; C2-18 (Alonso-González tr. 1, 11, 13); J33; C2-19 (Yu–Zhang–Zhou: chỉ tóm tắt trên Crossref)

Ba kết quả này đo ba đại lượng khác nhau: tới muộn so với giờ biểu, độ dao động của thời gian, và thông tin chờ trước
khi lên xe. Không kết quả nào cho một giá tiền của trần 30 giây trên phần dời giờ trả do điều phối, và cũng chưa có
chuẩn ngành cho loại trần này. Vì vậy khóa luận trình bày mỗi mức trần như một điểm trên đường đánh đổi, không coi 30
giây là ngưỡng khách cảm nhận được. Chương~\ref{chuong:phuongphap} đặc tả luật C trên nền các khái niệm của chương này.
% EV: D20; K52; J23 và J24 (lưu ý "hai cấu trúc khác nhau"); C2-23 (D03:260)
